\chapter{The language-model strand}
\label{chap:07-modeles-de-langue}

The previous chapters deal with programs that learn to act in a simulated world. This chapter
describes a smaller strand, devoted to the generative AI models that write text. It sets out the
vocabulary we use, the algorithm as we wrote it, the protocol, the runs and their results, then the
measurement trap these runs exposed. The full protocol is in \file{sim/llm/LLM\_PROTOCOL.md}. Its
twenty deviations are summarised in \cref{chap:09-ecarts}. Deviations~18 and~19, written
after all the results, declare defects found by a technical review and pick out in the
logs measurements that the first version of this chapter left aside. They change how the strand reads
without changing any published figure.

\section{Vocabulary}
\label{sec:llm-vocabulaire}

The vocabulary of RL has spread to generative AI models. Some RL algorithms are borrowed to refine
the answers of a language model, with no agent acting in a world. A language model is an AI model
that produces a text token by token, a token being a piece of a word. For a mathematics problem, a
program compares the final answer with the correct answer. The reward is 1 when the two agree.
The literature calls this RL with verifiable rewards (RLVR, \emph{reinforcement learning with
verifiable rewards}). We keep the acronym when we cite a paper that uses it, without adopting its
suggestion that the problem is the same as that of an agent in a world.

Three differences separate this setting from the classical RL of
\cref{chap:05-experiences-methode,chap:06-experiences-resultats}:
\begin{itemize}
  \item an \enquote{episode} shrinks to one answer: no state evolves under the actions, no
        reward is delayed over several decisions;
  \item the starting policy is a model already trained on huge corpora, which often answers
        correctly before any fine-tuning;
  \item the reward depends on the final text alone.
\end{itemize}
Fine-tuning from human feedback (RLHF, \emph{reinforcement learning from human feedback},
\cite{ouyang2022instructgpt}) is never called RL here: we refer to it by its acronym, as one
fine-tuning technique among others.

\section{Why this strand}
\label{sec:llm-pourquoi}

In 2025, \textcite{shao2025spurious} showed that a reward drawn at random, without checking whether
the answer is correct, improves Qwen models in mathematics, while Llama and OLMo models do not
improve. On the Qwen2.5-Math-7B model and the MATH-500 benchmark, the true reward brings a gain of
29.1 percentage points. A reward that praises a wrong answer fixed in advance brings 24.1, a reward
drawn at random 21.4. The authors trace this to a behaviour that training amplifies without creating
it: the model writes its reasoning as code, without running it, in 65\% of answers before training
and in more than 90\% after training with a misleading reward. They recommend validating a method on
several model families.

This result recalls two flaws of 2018. The best algorithm changes with the task (F13): here, the
model family plays the part of the task. Benchmarks that are too simple get learned (F16). The survey of \cref{chap:04-sondage-resultats}
measures the same risk in papers: sixteen of the 27 papers of the first 2026 draw that apply what
they call RL to large language models try one model family only (criterion X6). Our strand looks at
the phenomenon closely, at small scale, with two model families. The protocol announces its scope in
advance: an illustration on a laptop. It does not replicate \textcite{shao2025spurious}.

\section{The GRPO algorithm, as we wrote it}
\label{sec:grpo}

GRPO (\emph{Group Relative Policy Optimization}) comes from \textcite{shao2024deepseekmath}. It keeps
the clipped objective of PPO \parencite{schulman2017ppo} and replaces the critic network, which
estimates the value of a state, with a normalisation within a group of answers to the same problem.
We wrote it by hand, in PyTorch, in one file, \file{sim/llm/grpo.py}, to control every step. The file
had 159 lines for the runs described here. It has had 266 since the rerun of deviation~20 was added,
behind the option \texttt{-{}-c7}. Without this option, it repeats the earlier runs identically. No
fine-tuning library is involved: the file imports nothing but PyTorch and the Transformers library,
which loads the model and its tokeniser.

\subsection{The task and the reward}

The task we chose, \file{mul2}, asks for the product of two integers from 12 to 99. The file
enumerates the $88 \times 88 = 7{,}744$ ordered pairs, shuffles them with a generator whose random
seed is fixed (\texttt{random.Random(12345)}) and independent of the random seed of the run, then
sets the first 300 aside for testing. The other 7,444 serve for training. The prompt follows a fixed
template, \enquote{What is $a$ * $b$? Think step by step, then finish with the final answer as an
integer.}, passed through the model's chat template.

An answer is correct when the last integer it contains, thousands separators removed, equals the
product. The verdict requires no format, so that format is never rewarded. Let $c(o, q) \in \{0,1\}$
be this verdict for answer $o$ to problem $q$. The two conditions of the protocol define the reward
$r$ of each answer:
\begin{align}
  \text{true reward (\texttt{true})} &: \quad r = c(o, q),\\
  \text{reward drawn at random (\texttt{random})} &: \quad r \sim \mathrm{Bernoulli}(1/2),
  \text{ drawn without looking at } o.
\end{align}
The coin toss and the draw of each step's problems come from the same generator,
\texttt{random.Random(seed)}, which depends on the random seed of the run and nothing else.

\subsection{One training step}

Each step draws $B = 4$ training problems, without replacement within the step: a problem can come back
from one step to the next. For each problem $q_b$, the current policy $\pi_{\theta_{\text{old}}}$
produces $G = 8$ answers $o_{b,1}, \dots, o_{b,8}$, sampled at temperature 1, with neither \emph{top-p}
nor \emph{top-k}, limited to 160 tokens. For Llama, the sampled distribution is the model's own. For
Qwen, it is not (deviation~18, item~a). Qwen2.5's \file{generation\_config.json} file sets a repetition
penalty of 1.1, which \file{grpo.py} does not override. Before each draw, the logit $\ell$ of every token
already present in the text becomes $\ell/1.1$ if it is positive and $1.1\,\ell$ otherwise. Three
consequences follow.
\begin{itemize}
  \item Qwen's answers do not follow the policy whose probabilities $\log \pi_{\theta_{\text{old}}}$
        gives. Nothing in the computation corrects for this difference.
  \item The penalty hits the repeated digits of a computation.
  \item The comparison between the two families mixes the family with this decoding setting.
\end{itemize}

Each answer receives its reward $r_{b,i}$. The advantage of an answer compares its reward with those of
its group:
\begin{equation}
  \hat A_{b,i} = \frac{r_{b,i} - \bar r_b}{s_b + 10^{-4}}, \qquad
  \bar r_b = \frac1G \sum_{j=1}^{G} r_{b,j}, \qquad
  s_b = \Bigl(\frac{1}{G-1}\sum_{j=1}^{G} (r_{b,j} - \bar r_b)^2\Bigr)^{1/2}.
\end{equation}
The standard deviation $s_b$ follows the PyTorch convention, with the factor $G-1$. A group in which
all eight answers receive the same reward gives each of them a zero advantage and contributes no
gradient.

For token $t$ of answer $o_{b,i}$, the probability ratio compares the policy being updated with the
policy that produced the answer:
\begin{equation}
  \rho_{b,i,t}(\theta) = \exp\bigl(\log \pi_\theta(o_{b,i,t} \mid q_b, o_{b,i,<t})
    - \log \pi_{\theta_{\text{old}}}(o_{b,i,t} \mid q_b, o_{b,i,<t})\bigr).
\end{equation}
The loss we minimise is the negative of the clipped objective, averaged over tokens within each
answer, then over answers within the step:
\begin{equation}
  \mathcal L(\theta) = -\frac{1}{BG} \sum_{b=1}^{B} \sum_{i=1}^{G} \frac{1}{|o_{b,i}|}
  \sum_{t=1}^{|o_{b,i}|} \min\Bigl(\rho_{b,i,t}(\theta)\,\hat A_{b,i},\;
  \operatorname{clip}\bigl(\rho_{b,i,t}(\theta), 1-\varepsilon, 1+\varepsilon\bigr)\,\hat A_{b,i}\Bigr),
  \qquad \varepsilon = 0.2.
  \label{eq:grpo}
\end{equation}
The length $|o_{b,i}|$ counts the tokens of the answer that are not padding (the mask
\texttt{target != pad\_token\_id} in the function \texttt{token\_logps}). In Qwen, the padding token is
also the stop token \texttt{<|endoftext|>}: the mask thus removes from the gradient the token that ends
the answer (deviation~18, item~d). The log-probabilities of the old policy, $\log \pi_{\theta_{\text{old}}}$,
are computed once, without gradient, before the first update.

Each batch of 32 answers serves for two update passes (\texttt{EPOCHS = 2}). In each pass, the
gradient accumulates over sub-batches of eight answers (two for Llama, setting \texttt{MICRO}), which
gives the same gradient with less memory. The gradient norm is then clipped at 1. The Adam
optimiser \parencite{kingma2015adam}, with PyTorch's default settings ($\beta_1 = 0.9$,
$\beta_2 = 0.999$, $\varepsilon = 10^{-8}$), applies a learning rate of $3 \times 10^{-6}$. The weights
stay in 32-bit floating point.

\begin{notebox}
\textbf{How our GRPO differs from that of \textcite{shao2024deepseekmath}.} Our loss has no
Kullback-Leibler divergence term with respect to the starting model, as the protocol
announced. The group holds 8 answers per problem and the step 4 problems, far fewer than in the
reference papers. The 160-token limit is very short for written reasoning. The learning rate was not
chosen by the planned pilot, which was stopped at step 10 (deviations 4 and 6). These choices are
adequate for an illustration. They rule out reading our figures as a measurement of GRPO at the
scale of the papers.
\end{notebox}

\subsection{Safeguards against a corrupted computation}

The first runs on the Mac's graphics processor (GPU, through PyTorch's MPS interface) crashed on
non-finite probabilities. Deviation~10 of the protocol documents the cause: macOS interrupts GPU work
that hinders interactive use of the machine. Such a computation sometimes comes back without an
error, with corrupted values. The same step, redone on the central processor (CPU), gives a gradient
norm of 14.4 before clipping, with no non-finite value. Since that deviation, \file{grpo.py} applies
four rules:
\begin{enumerate}
  \item an update whose gradient norm before clipping is non-finite or exceeds 1,000 is recomputed,
        up to five times in a row, after which the run stops (\texttt{GRAD\_SANE},
        \texttt{RETRIES});
  \item a generation that raises a GPU error is redone, up to five times;
  \item the run stops if the log-probabilities of the old policy or the weights become non-finite;
  \item the log counts the recomputations at each step (field \texttt{retries}).
\end{enumerate}
The threshold of 1,000 was set from the norm of 14.4 measured on CPU. A gradient that is corrupted but
finite, below this threshold, then clipped at 1, would pass without a trace. The protocol records
this as a limit, in particular for the first three runs, which ran before these rules existed.

\subsection{The measurement}

Evaluation uses the test set and combines two modes. Greedy accuracy takes one answer with no
randomness, the most probable token at each position, for each of the first 150 test problems. Sampled
accuracy draws 8 answers at temperature 1 for each of the first 50 test problems. With $c_{jk}$ the
verdict on answer $k$ to problem $j$ and $\tau_{jk}$ the indicator of a cut-off answer:
\begin{equation}
  \text{pass@1} = \frac{1}{50}\sum_{j=1}^{50} \frac18 \sum_{k=1}^{8} c_{jk}, \qquad
  \text{pass@8} = \frac{1}{50}\sum_{j=1}^{50} \max_{k} c_{jk}, \qquad
  \text{cut-off share} = \frac{1}{400}\sum_{j,k} \tau_{jk}.
\end{equation}
An answer is cut off when its 160 tokens contain neither an end token nor a padding token (function
\texttt{truncated}). This function knows only the main end token: it ignores Llama's two extra stop
tokens, 128001 and 128008 (deviation~18, item~d). In Qwen, the repetition penalty also applies to the
greedy answer. Here pass@8 is the observed frequency over exactly eight draws, without the unbiased
estimator used when more than $k$ answers are drawn. Evaluation takes place at step 0, step 50 and
step 100 (\texttt{EVAL\_EVERY = 50}). Every ten steps, the log also records the share of correct and of
cut-off answers in the training batch, the loss, the gradient norm, the recomputations and the elapsed
time. Each run writes a JSONL file, \file{sim/llm/results/<model>/mul2/<reward>/seed<seed>.jsonl}. The log
keeps the four measures of each evaluation, but no verdict per problem. Evaluation draws its answers from
the same random generator as training: the evaluation at step 50 changes the draws of the following steps
(deviation~18, item~d).

The main measure is sampled pass@1, the quantity GRPO trains (deviation 3). The first version of this
chapter left greedy accuracy and pass@8 in the logs. Deviation~19 reports them, with a fifth measure,
the accuracy among finished answers:
\begin{equation}
  \text{finished accuracy} = \frac{\text{pass@1}}{1 - \text{cut-off share}},
\end{equation}
which assumes that a cut-off answer is never correct (0 out of 23 in the diagnosis of deviation~5,
\cref{sec:llm-coupure}). It separates what pass@1 owes to accuracy from what it owes to brevity.
Pass@8 answers the question of \textcite{yue2025rlvr}: does fine-tuning widen what the model can do, or
does it concentrate probability on answers the model already found within several draws?

\section{The protocol and its runs}
\label{sec:llm-protocole}

\paragraph{Models} Two families, two small open models: Qwen2.5 with 0.5 billion parameters,
\emph{Instruct} version \parencite{qwen2025qwen25}, then Llama 3.2 with 1 billion, \emph{Instruct}
version, in the copy published by \texttt{unsloth} \parencite{meta2024llama32,grattafiori2024llama3}.
Since deviation 15 (c), \file{grpo.py} pins the exact revision of the downloaded files
(\texttt{7ae5576\dots} for Qwen, \texttt{5a8abab\dots} for Llama), the only revisions present in the
machine's cache, downloaded before the first run.

\paragraph{Pilot} A pilot without training was to fix the task: the one on which both models lie
between 15\% and 70\% accuracy. The protocol wanted it on problems outside the test set. It ran on the
test set (deviation~18, item~b): the option \texttt{-{}-pilot} calls the evaluation function, which
reads the first 150 and 50 test problems. The choice of \file{mul2} and that of the 160-token limit were
thus made while looking at the test set. No training run saw the test set. In greedy decoding, no task
places both models in this range. The sum of three numbers (\file{add3}) is too easy for Llama. The
multiplication \file{mul2} gives 7\% (Qwen) and 73\% (Llama) in greedy decoding, 11\% and 23\% in
sampled pass@1 (deviation 3). Deviation~3 retains \file{mul2} all the same. The same pilot showed that 64 tokens cut off the answers of Qwen, which
writes 100 to 160, for an apparent accuracy of 0.7\%. The limit went up to 160 tokens, after reading
real outputs, before any training (deviation 1).

\paragraph{Plan} Two models, two rewards, 100 steps per run, a learning rate of $3 \times 10^{-6}$
for all conditions. The protocol planned 200 steps and five random seeds. Deviation 6 reduced the
plan to 100 steps and three random seeds (0 to 2), before any retained run, to fit the available
compute. Deviation 14, written before any new run, restored random seeds 3 and
4, with the same code and the same settings. The starting point of each run is the base model,
measured with the random seed of the run.

\paragraph{Statistics} For each random seed $s$, the gain is $g_s = \text{pass@1}_{100} -
\text{pass@1}_{0}$. Each configuration (model, reward) is summarised by the mean gain and a 95\%
percentile bootstrap interval: 10,000 resamples of the gains with replacement, generator
\texttt{random.Random(0)}, bounds at ranks 250 and 9,750 of the sorted means
(\file{sim/llm/summary.py}). The details of the bootstrap were fixed after the first results, with
no choice among several variants (deviation 13). The interval is very rough: three random seeds
allow only ten distinct resamples, five random seeds 126.

\paragraph{Added statistics} Deviation~19, written before the script \file{sim/llm/logs\_stats.py},
adds three tools for the five measures. The technical review that preceded it had displayed some of these
values: the reading rule does not depend on their sign.
\begin{itemize}
  \item Within each configuration, beside the percentile interval, the Student interval
    $\bar g \pm t_{n-1;\,0.975}\, s_g/\sqrt n$, where $\bar g$ and $s_g$ are the mean and the standard
    deviation of the $n$ gains. With five random seeds, the percentile interval is only about 60\% as wide
    as the Student interval and covers too little. The exact sign test compares the gains with the $2^n$
    possible sets of signs.
  \item Between the two rewards, the question the protocol itself asks: the difference of the mean gains,
    its Welch interval, with the Welch-Satterthwaite degrees of freedom, then an exact permutation test over
    every split of the gains between the two rewards, 252 for Qwen and 126 for Llama. Holm's correction
    applies to the two models for pass@1, the main contrast, as it does for the finished accuracy. It then
    applies to the other eight contrasts.
  \item A rule: a gain or a contrast is \emph{established} when its Student or Welch interval excludes 0
    and its exact $p$, corrected by Holm, falls below 0.05.
\end{itemize}
This rule has a consequence the protocol had not seen. With five random seeds, the exact sign test never
goes below $p = 2/2^5 = 0.0625$, reached when all five gains have the same sign. With four, it does not go
below $2/2^4 = 0.125$. No gain within a configuration can therefore be established with five random seeds,
even when all of them point the same way. More answers per problem would change nothing. At step 0, the
standard deviation of pass@1 between random seeds is 0.7~points for Qwen and 1.1~points for Llama, below
the binomial standard error of 400 answers (1.5 and 2.1~points). Measurement noise is therefore small next
to the spread of the gains under the true reward, whose standard deviation is about 24~points for Qwen and
10 for Llama. The precision of the mean gain depends on the number of random seeds far more than on the
number of answers.

\paragraph{A run that fails twice} The Llama run with the random reward and random seed 4 stopped
twice in a row on the first safeguard, at its first update. Deviation 16 counts this second
failure as a result: the configuration is summarised over its four finished runs, with the failure
stated next to them. A sensitivity analysis also replaces the missing gain with the smallest, then
with the largest gain of the nineteen finished runs.

\paragraph{Hardware} All runs ran on the project's Mac
(\cref{chap:08-logiciel,chap:D-environnements}). The first three Qwen runs ran on the GPU, on an
idle machine. After the crashes of deviation 8 and the diagnosis of deviation 10,
the remaining Qwen runs moved to the CPU, with two threads (\file{cpu.sh}, deviation 11). The Llama
runs, about two and a half times slower than Qwen on CPU, stayed on the GPU, on an idle machine only
(\file{night.sh}), with the safeguards. \Cref{tab:grpo} gives the device of each run.

\section{Results}
\label{sec:llm-resultats}

\gentable{grpo}

\Cref{tab:grpo} gives, for each of the twenty runs, pass@1 at the start and after 100 steps, the
device and the share of cut-off answers at the end. \Cref{fig:grpo-graines} shows each run as an
arrow, \cref{fig:grpo-courbes} the three evaluation points of each run.

\fig{grpo-graines}{Pass@1 at the start and after a hundred training steps, for each model, each
reward and each of the five random seeds. A filled circle marks a run on GPU, an empty square a run
on CPU, a cross the run that stopped twice. With the true reward, the same program gives very
different results from one random seed to another. With the coin toss, success falls to zero or
nearly so.}

\fig{grpo-courbes}{Test pass@1 at steps 0, 50 and 100 of each run, per model, reward and random
seed. Evaluation takes place only every fifty steps, so each curve joins three points.}

\subsection{The reward drawn at random}

Under the reward drawn at random, success falls to 1.25\% or less in all nine finished runs. The mean gain is $-10$
points on Qwen (from $-11$ to $-9$) and $-24$ points on Llama (from $-24$ to $-23$). Replacing the
missing Llama gain with the smallest gain of the nineteen finished runs changes little ($-24$, from
$-24$ to $-23$). Replacing it with the largest, $+51$, lifts the mean to $-9$, with an interval from
$-24$ to $+21$. In our hands, at this scale, the reward drawn at random destroys the starting
competence of both families. It does not reproduce the gain that \textcite{shao2025spurious} observe
on Qwen2.5-Math-7B. Our conditions differ from theirs in almost every respect (size and
specialisation of the model, task, answer length, number of steps, absence of a KL term), so this
disagreement does not refute their result.

\subsection{The true reward, in two stages}

With three random seeds, Qwen starts at about 10\% and ends at 36\%, 61\% and 13\%. The mean gain is
$+26$ points, with an interval from $+2$ to $+51$. On Llama, success starts between 23\% and 25\% and
ends at 20\%, 48\% and 25\%, a mean gain of $+7$ points, from $-4$ to $+22$. The same program gives
three very different results: this is flaw F11, same code with another random seed, in miniature.

With two more random seeds, the new Qwen runs end at 13\% and 0\%. Over five random seeds, the mean
gain falls to $+14$ points, from $-3$ to $+34$: the interval now contains zero. The gain we read
on Qwen with three random seeds no longer holds. On Llama, the two new runs end at 27\% and 33\%. The
mean gain is $+6$ points, from $-0.4$ to $+15$: zero stays inside the interval. The Student intervals
of deviation~19 are wider: from $-16$ to $+45$ for Qwen, from $-6$ to $+19$ for Llama, with a sign-test
$p$ of 0.31 for both. With five random seeds, we cannot say that GRPO improves either model on this
task, under these conditions.

\subsection{What pass@1 hid}

\Cref{tab:grpo-mesures} gathers the five measures of each configuration, \cref{fig:grpo-mesures} their
gains seed by seed. Under the true reward, Llama gains 6~points of pass@1 while its greedy accuracy drops
in all five runs, from 72.1\% to 29.1\% on average ($-43$~points, Student $[-59, -27]$). Its pass@8 drops
too, from 75.2\% to 44.8\% ($-30$~points, $[-44, -17]$). The model does not compute any better. The
distribution it samples at temperature 1 narrows: a random draw is right a little more often, while the
most probable answer and the best of eight are right less often. In Qwen, pass@8 falls in four runs out of
five, from 51.6\% to 34.4\%, with a Student interval that contains zero. To the question of
\textcite{yue2025rlvr}, our runs answer that pass@8 rises on average in no configuration. Under the rule of
deviation~19, Llama's two drops are not established, since the sign-test $p$ cannot fall below 0.0625,
while their Student intervals exclude zero. Here the rule therefore declines to establish drops that the
Student interval, taken alone, would flag.

\gentable{grpo-mesures}

\fig{grpo-mesures}{Gains per random seed on three measures, between step 0 and step 100, for Qwen2.5 with
0.5 billion parameters and Llama 3.2 with 1 billion, under the true reward and under the coin-toss reward.
One dot per random seed, the mean, the Student interval (thick line) and the protocol's percentile
interval (thin line). Under the true reward, Llama's pass@1 rises by 6~points while its greedy accuracy
loses 43~points and its pass@8 30: the sampled distribution narrows, the model does not multiply any
better. With five random seeds, the percentile interval is much narrower than the Student interval.}

\subsection{True reward against coin toss}

The protocol asked a question that the first version of this chapter did not test: does the true reward
do better than the reward drawn at random? \Cref{tab:grpo-contraste} answers through the difference of the
mean gains. For Llama, the true reward does better in pass@1, by $+30$~points $[+18, +43]$ (exact
permutation $p = 0.008$, Holm 0.016): this is established. For Qwen, the difference points the same way,
$+25$~points $[-6, +55]$ ($p = 0.040$, Holm 0.040), but its interval contains zero: it is not
established. None of the eight secondary contrasts passes Holm. Why Llama's contrast exists matters as
much as its sign. The true reward does not make Llama progress: it degrades Llama less than the reward
drawn at random, which destroys it (pass@1 from 24\% to 0.3\%, greedy accuracy from 72\% to 0.8\%).

\gentable{grpo-contraste}

\subsection{Hardware and random seed, mixed up}
\label{sec:llm-materiel}

The Qwen gain with the true reward comes from the two runs made on GPU, random seeds 0 and 1: $+26$
and $+51$ points. The three true-reward runs that ran on CPU (random seeds 2 to 4) have a mean gain
of $-2$ points, from $-11$ to $+4$. The four CPU runs with the random reward (random seeds 1 to 4)
have a mean gain of $-11$ points, from $-11$ to $-10$. On CPU alone, the true reward brings Qwen no
detectable gain. Data: \file{sim/llm/results/gains\_5.json} and \file{extra\_5.json}.

These figures do not separate two causes. The hardware changed with the random seed: each Qwen
random seed with the true reward ran on one device only. Two explanations remain open. In the
first, chance in the random seeds explains everything: the two favourable random seeds happened to
fall on the GPU. In the second, the device changes the course of training, through the rounding of the
computations or through GPU computations that were corrupted but finite and that the first runs
would not have detected. Five runs do not decide between these explanations.

The logs establish two facts about the device (deviation~18, item~c). Pass@1 at step 0 depends on the
base model, the random seed and the chance of the draws, since the reward comes into play only after
this first evaluation. Greedy accuracy at step 0, which draws nothing at random, is \num{7.33}\% in all
ten Qwen runs, on CPU as on GPU. First fact, then: the device does not change the starting model's most
probable answer. The differences in pass@1 at step 0 between Qwen runs with the same random seed follow from
the random streams, which differ between CPU and GPU for one seed. For Qwen with random seed 1, for instance,
the true-reward run, on GPU, starts from \num{10.25}\%, the random-reward run, on CPU, from \num{11.75}\%.
The second fact, the direct sign that the MPS computation is not deterministic, comes from Llama. For random
seed 3, both runs ran on GPU, with the same model and the same random stream. They nevertheless start from a
pass@1 of 25\% and \num{24.25}\%. Their greedy accuracy is \num{70.7}\% on one side, \num{66.7}\% on the
other.
Deviation~7 had already observed it (\cref{tab:grpo}).

The logs also show groups that yield no gradient. Under the true reward, many groups have eight answers
all wrong, or all right, hence a zero advantage for each. The gradient norm is zero in 5, 4 and 7 of the
10 logged steps for the Qwen runs with random seeds 2, 3 and 4 (random seeds 0 and 1 did not log it). It is
zero in 3, 0, 3, 3 and 2 steps out of 10 for the five Llama runs. The Qwen run with random seed 4 stops
learning from step 50, at 0\% pass@1, pass@8 and greedy accuracy. Under the reward drawn at random, the
coin almost always gives mixed rewards: no logged step has a zero gradient.

\paragraph{A run outside the protocol} Deviation 16 (e) planned to run the Llama run with the random
reward and random seed 4 once on CPU, after all other computations, to learn whether its double
failure was due to the GPU. The run completed with no recomputation. Its pass@1
goes from \num{23.25}\% at step 0 to \num{0.25}\% at step 100
(\file{sim/llm/results/llama1\_cpu/mul2/random/seed4.jsonl}). The result stays out of every
statistic, as the deviation had set. It is consistent with a failure due to the GPU without proving
it, since one run does not separate the device from chance. On CPU, this tenth random-reward
run behaves like the other nine: success collapses.

\section{The cutoff trap}
\label{sec:llm-coupure}

\subsection{The diagnosis}

The incident of deviation 4 exposed this trap. During the learning-rate pilot, the share of correct
answers in the training batch was 0.91 at step 10, while the test pass@1 of the base model was 0.11.
The diagnosis of deviation 5 (\file{sim/llm/diag/replay\_batch.py}) replays, without training, the
exact batch of step 10 of the pilot: the four multiplications $48 \times 49$, $14 \times 36$,
$21 \times 50$ and $45 \times 17$, eight answers each, with the starting Qwen model. The model gives
5 correct answers out of 32, in line with the test. Yet 23 of the 32 answers reach the 160-token
limit before giving a result. None of these cut-off answers is correct. Among the 9 complete
answers, 5 are correct (\cref{fig:grpo-coupure}).

\fig{grpo-coupure}{The answers of the Qwen2.5 model with 0.5 billion parameters before any training,
eight per problem, on the four multiplications of the tenth step of the pilot. A complete correct
answer is green, a complete wrong answer grey, a cut-off answer red. Most answers reach the 160-token
limit before giving a result and are scored as wrong.}


The 0.91 of step 10 was no counting error. The simplest explanation is that the true reward first
teaches the model to finish before the cutoff. Yet when the starting model replays this batch, 5 complete
answers out of 9 are correct, against no cut-off answer. The reward drawn at random has no such lever, since it never looks at the answer. The
protocol kept the 160 tokens and added the measurement of the cut-off share, in training and in
test, before any retained run.

\subsection{What the twenty runs show}

At the start, the share of cut-off answers is high for Qwen and low for Llama: from 71\% to 77\% for
Qwen, 3\% at most for Llama, according to the evaluation logs at step 0. With the true reward, over
the ten runs, it ends between 0\% and 29\%. Under the reward drawn at random, it ends above one half in
six runs out of nine, between 2\% and 16\% in the other three. Two Llama runs with the random reward end with
99\% and 100\% of answers cut off. A third cuts off only 2\% and no longer gives any correct
answer: the reward drawn at random does more than make the answers overrun the limit.

A cut-off answer counts as wrong. A model that learns to answer more briefly therefore gains pass@1
points without computing any better. The first version of this chapter could not put a figure on this
part. The finished accuracy of deviation~19 does (\cref{tab:grpo-mesures}).
\begin{itemize}
  \item In Qwen, under the true reward, the cut-off share falls from 74.9\% to 6.0\% in all five runs. The
    accuracy among finished answers \emph{falls}, from 41.2\% to 27.5\% ($-14$~points, Student
    $[-46, +19]$, two runs out of five up). The end of cut-off answers accounts for the whole pass@1 gain
    of Qwen. The model learns to finish within 160 tokens. It does not multiply any better.
  \item In Llama, the length barely moves: the cut-off share goes from 2.0\% to 2.7\%. The finished
    accuracy follows pass@1, $+7$~points, Student $[-5, +19]$.
\end{itemize}
Within the Qwen configuration with the true reward, the Pearson correlation between the change in the
cut-off share and the change in pass@1 is $+0.23$ over five runs. On so few points, nothing suggests that
shortening further brings a further gain.
Over both rewards together, the Spearman correlation is $-0.63$ for Qwen and $-0.83$ for Llama, which
mostly reflects the effect of the reward (\file{sim/llm/results/logs\_stats\_5.json}, key
\texttt{length}).

\begin{keybox}
\textbf{Lesson for a language-model benchmark.} A length limit turns a measure of accuracy into a
mixed measure of accuracy and concision. A paper that reports a pass@1 gain is read together with the
length limit, the share of cut-off answers before and after training and the accuracy of
the complete answers. Without these three numbers, the reader does not know what the gain measures.
\end{keybox}

\section{What this strand establishes and what it does not}
\label{sec:llm-portee}

This strand establishes four facts, for two small models, one multiplication task and 100 steps.
\begin{itemize}
  \item With the true reward, the dispersion between random seeds is as large as the mean pass@1 gain, as in
    the classical RL of 2018. No pass@1 gain is established.
  \item A reward drawn at random makes the success of both families collapse, in every run and on all
    four accuracy measures where they are defined. One Llama run (random seed~1) finishes no answer, so
    accuracy among finished answers does not exist there and is computed on random seeds 0, 2 and 3 only.
    This collapse gives Llama its only established contrast between the two rewards.
  \item In Llama, the true reward narrows the sampled distribution: pass@1 rises a little, while greedy
    accuracy and pass@8 drop in all five runs.
  \item In Qwen, the whole pass@1 gain is a gain in brevity, since the finished accuracy does not rise.
\end{itemize}
Two declared defects weigh on these facts: Qwen's repetition penalty and a pilot that read the test set
(deviation~18). The five-seed protocol cannot establish any gain within a configuration. A rerun is
written in advance to remove these defects (deviation~20): repetition penalty at 1, pilot kept off the
test set, ten random seeds per configuration on the GPU alone, verdicts kept per problem and a greedy
answer at 512 tokens, which the length limit does not touch. It has been launched, is waiting
for the GPU to become free and has produced no result yet. Its rule is fixed: if the main contrast is
established for neither model, this strand moves to an appendix.

The strand establishes neither the effect of GRPO at the scale of the papers nor the absence of the
effect of \textcite{shao2025spurious} on larger models or on models specialised in mathematics. It does
not separate the effect of the hardware from that of the random seed for Qwen. The limits that follow
are taken up in \cref{chap:10-limites}.
